\documentclass[11pt, a5paper, openright]{book}

% This is a transcription done by Mairi Dulaney of Seattle, WA

% chktex-file 18

\usepackage{caption}
\usepackage[english]{babel}
\usepackage[T1]{fontenc}
\usepackage[]{geometry}
\usepackage{textcomp}
\usepackage{gensymb}
\usepackage{graphicx}
\graphicspath{{images/}}
\usepackage[utf8]{inputenc}
\usepackage{makeidx}
\usepackage[round,sort,comma]{natbib}
\usepackage{rotating}
\usepackage{titlesec}

\usepackage[activate={true,nocompatibility},final,
  tracking=true,kerning=true,factor=1100,stretch=10,shrink=10]{microtype}

\geometry{a5paper}
\nonfrenchspacing

\titleformat{\part}[display]{\centering\normalfont\Huge\bfseries}{}{0pt}{}

\titleformat{\chapter}[display]
{\bfseries\Large}
{\centering
  \textsc{Chapter} \thechapter.}
{0.5ex}
{
  \vspace{1ex}
  \centering
  \scshape
}
[\vspace{-1ex}]

\titleformat{\section}[display]{\centering\bf}{\enspace}{0em}{}

%\overfullrule=2cm

\renewcommand{\thefootnote}{\fnsymbol{footnote}}


% Some macros to make indexit easier
\newcommand{\barge}[1]{\textit{#1}\index{#1,~barge}}
\newcommand{\steamer}[1]{\textit{#1}\index{#1,~steamer}}

\makeindex


\begin{document}

\renewcommand{\thechapter}{\Roman{chapter}}
\title{Riverboating in Lower Carolina}
\author{F. Roy Johnson}
\date{1977}
\bibliographystyle{alpha}
\maketitle


\part{I}

\chapter{By Paddle, Ore, and Pole}

\section{Early Boats of Burden}

\textsc{Although} the dugout canoe, or log boat, has disappeared from
the rivers, creeks, and other waters of North Carolina's southeastern
counties, the watersheds of the Cape Fear, Black, and Northeast
rivers, it must be remembered as the common and more important vehicle
of travel and commerce of the early settlers of the region.  And to a
diminishing extent this versatile vehicle continued in common service
far into the nineteenth century.  As late as the 1950's a few
serviceable old moss cloaked log boats reposed in the South River
wilds of Bladen and Sampson counties and at remote landings of the
upper Waccamaw River in Columbus County.\par

The numerous waterways of the southeastern area --- the sounds, the
rivers, and the creeks --- were well suited for small boats.  These
avenues grew into the great channels of commerce and retained their
commanding position for nearly two centuries, until the railroads
became common and transportation needs moved westward.\par

Although much of the early travel was overland by horseback and by
foot on trails little better than Indian paths, virtually all the
commerce was by water.  However, by the Revolution, overland traffic
had increased; for the industrious settlers had built roads,
constructed causeways across swamps, bridged streams, and established
ferries.  This lessened little the water freightage.  And up and down
the waterways there plied a great variety of craft loaded with
commodities.\par

Plantations, trading places, villages, and towns sprang up along the
streams; and the growth of the towns was determined to a large extent
by the volume of this commerce.  Chief of the towns were Fayetteville
at the head of navigation on the Cape Fear River, which funneled goods
down stream from the back country and received imported commodities,
and Wilmington, which exported and imported goods for the watersheds
of the Cape Fear River and her tributaries.\par

Not so fortunate was the town of Lisbon at the head of the Black River
in Sampson County.  Its merchants prospered and the town grew and
flourished on pole boat traffic.  But after steamboats came it went
neglected because of the river's shallows; and it dried up completely
when bypassed by the railroads.\par

\section{The Canoe and Periauger}

The canoe's larger companion, the periauger, was a boat of burden
which served in the early commerce of the southeastern counties.  And
as the volume of commerce increased it was followed by the larger flat
boat, or pole boat.\par

Writing in 1737, a time when settlers were pouring into the
southeastern area, Dr.\ John Brickell told of the building and use of
the canoe and the periauger.  The canoe, eh said, was made out of one
piece of timber and commonly of cypress, the durable giant of the
coastal swamp forests.  A log of proper thickness and length was cut
from the felled tree, and craftsmen shaped the log like a boat and
then hollowed it out.  Sometimes the boat was made wider by splitting
it down the middle and installing boards in its bottom; and this boat
was called the periauger.  Then the canoe or periauger was fitted with
masts for sailing or oars and paddles, ``according to \ldots size and
bigness.''\par

Some of the larger of these craft were ``capable of carrying forty or
fifty Barrels of Pitch and Tar.''  The planters also found the
periauger useful as ferries and for transporting goods from one
plantation to another.  The design of this great log boat made it so
efficient that ``no Vessel of the same Burthen made after the European
manner'' could outsail one.\citep[260-261]{brickellj}\par

William Robinson, a Sampson County planter who lived at the Delta on
the Black River about sixty miles from Wilmington by water, was using
a large cypress log boat in regular commerce with that seaport town in
the late eighteenth century.  He manned her with six black oarsmen.
And late in the nineteenth century William's great-grandson Scott
Robinson was running a turpentine distillery at the Delta and
freighting naval stores to Wilmington on a fifty foot by twelve foot
flat poled by four blacks.  \citep{shawg}\par

The canoe was made in many sizes.  The smaller might carry no more
than two or three men while the larger could carry two or three
horses.  They seem to have continued quite numerous throughout the
eighteenth century.  During the spring of 1762 a group of Moravians
proceeding up the Cape Fear River watched their boats by night for
fear of Negroes who ``go about in small canoes, stealing where they
can.'' \citep[I,~261]{friesal}.  And in 1775 Janet Schaw, a Scotch
woman sojourning in North Carolina, saw more than one hundred people of
the lower class in canoes come to a funeral on the Northeast
River.\citep[171]{schawj}\par

The larger boats used sails to great advantage, except on narrow and
croked streams where oars and poles were needed.  Tides favored all
kinds of craft upon the lower rivers, and freshet waters often aided
navigation upon their headwaters.\par

\section{Eighteenth Century Row Boats}

Late in the eighteenth century large \textit{row boats} were the
carriers of produce on the Cape Fear River between Cross Creek, near
present Fayetteville, and the Wilmington area \citep[280-281]{schawj}.
The planter near Wilmington manned such boats with slaves.  For travel
one such planter had ``a very fine boat with an awning to prevent the
heat and six stout Negroes in neat uniforms to row her down, which with
the assistance of the tide was performed in a very short time''
\citep[177]{schawj}.\par

As the volume of commerce increased, so did the size of the boats.
This led to the construction of the flat boat, or \textit{flat}, so
called because of its flat bottom and squared sides, which was built
of boards.  Being of shallow draught it could be used on small
streams.\par

Navigation of the streams often was beset with problems.  Hugh
Meredith, who about 1730 had come to North Carolina from Pennsylvania
to settle, wrote that the channel of the Northeast Cape Fear River was
impeded by the ``multitude of Logs that lie in it, part of them fast in
the Sand, with great Snags or Limbs, and sometimes either End the
Middle quite above, or but little beneath the Surface; and in some
places we saw whole Heaps jumbled together, almost from Side to Side,
and so firm that they are immovable, being sound, heavy, fast, and
deep in the Sand.'' (Meredith, Hugh, 21-22).\par % chktex 8

Soon after the Revolution there was a burst of enthusiasm for
improvement of waterways for navigation.  Some progress was made in
making passable the channels of the main streams.  And by the opening
of the nineteenth century virtually every little stream was cleared of
obstructions so as to accomodate small boats and flats, giving many
farmers located off the large streams access to water
transportation.\par

\section{The Pole Boat}

South (Black) River, a tributary of the Black River, which rises in
Harnett County, was an important avenue of commerce for farmers of
eastern Cumberland County.  They ran their rafts and pole boats upon
it until late in the nineteenth century when railroads provided an
alternate means of transportation.\par

Caleb C. Bullard, born in 1859 in Beaverdam Township, recalled that
naval stores were carried down from this area to Wilmington on flat
boats.  John A. Oates in his \textit{The Story of Fayetteville} quotes
Bullard as saying, ``There was no trouble in getting down the river,
except to keep the flat boat guided; there was no steam power on South
River.  It took usually about four days for the trip down and six days
for the trip back, about 100 miles each way.  In Wilmington, they
would secure, usually, a tug boat to bring them to the mouth of South
River and from there up they would use poles twenty feet long or more,
to push the flat boat on up the river.  They would have six or eight
men to do this.''\par

Naval stores were the chief source of income for the Cumberland
farmers until some time after the Civil War when they began to raise
cotton.  In the depressed economy they found little inducement to grow
more than enough food to feed their families and livestock
\citep[432]{oatesja}.\par

\section{The Large Flat of Tidewater}

In tidewater no other power than the tide was needed to move the large
flat.  It was tied up against the tide, and when the tide turned it
floated like a raft with it.  Large flats, some carrying as much as 100
to 200 barrels of turpentine, tar, and pitch, were in use upon the
lower Cape Fear before the Revolution \citep[184-185]{schawj}, but
smaller craft were common above tidewater until the coming of the
steamboat which could tow them upstream against the current.\par

Some idea of flat boat sizes early in the nineteenth century may be
obtained from an advertisement in the November 21, 1816
\textit{Fayetteville American}.  An advertiser was seeking to let a
contract for construction of four flat boats, two 48 feet by ten feet
and two 30 feet by six feet.  They were to be made of heavy timber and
sides banded with strap iron \citep[11-21-16]{fa}.\par

Midway through the nineteenth century tide powered flats were carrying
300 to 500 barrels of turpentine or an equal amount of other
cargo.\footnote{The Wilmington Journal of 1851, under ``Marine
  Intelligence'', provides a record of pole boat activity.}  For
example, flat boat
\textit{J.~L.~Cassidy}\index{J.~L.~Cassidy, flatboat} from Lyons
Landing on the lower Black River arrived in Wilmington on April 11,
1851 with 326 barrels of rosin, 64 barrels of tar, and 64 barrels
of turpentine.  By comparison, on January 31 the flat boat
\textit{Sopha}, with Man William in charge, from Beatty's Bridge on
the Black River 50 miles above Wilmington and above tidewater, arrived
with 38 bales of hay, five bales of fodder, and four sheafs of oats.
It's cargo was comparable to about ten barrels of turpentine.\par

Meanwhile, the large tide powered flat boat was hauling much of the
cargo on the lower Cape Fear, Black, and Northeast rivers.  White
Hall, on the Cape Fear about 18 miles below present Elizabethtown and
54 miles above Wilmington became an important receiving and trading
center.  With some assistance from poles the great flats were able to
make their way this far up with no great difficulty.  Those going
higher up mostly were towed by steamboats.  Some of these, especially
those owned by the steamboat companies, were towed back down stream;
and these were called \textit{tow boats}.\par

Often the large flat boat was given a name which honored its owner or
some other person and put in charge of a trusted slave.  In 1851 one
finds Man Dick in charge of \textit{Uncle Sam}; Man Sandy, the
\textit{David Lewis}, Man Bob, the \textit{Stevenson}; Man Wilson, the
\textit{General Taylor}; Man Jack, the \textit{Dried Apple}; Man
Daniel, the \textit{James Ellis}; Man Dick, the
\textit{J.~L.~Cassidy}.  Other big flats were the \textit{David Reed},
the \textit{Jackson}, the \textit{General Cass}, \textit{Macmillan's
  Boat}, and \textit{Robinson's Boat}.\par

\section{The Steam Flat}

The steam flat followed the pole and row boat and small flat into the
narrow and shallow headwaters of the Black and Northeast rivers and
large creeks to compete with rafts in bringing down naval stores from
the developing country.  Although some of the early steamboats on the
Cape Fear River were little more than steam-powered paddle-wheel
flats, the steam flat as a distinctively different kind of craft did
not appear until after the Civil War when the entire southeastern area
of North Carolina was faced with the need for economical
transportation.\par

In 1886 ten of these vehicles of burden were transporting naval stores
and other commodities from the back country of Bladen, Brunswick,
Pender, Sampson, and Onslow counties to Wilmington.  Two of these were
operated on the Great Coharie, a shallow tributary of the Black River
up to 100 miles from Wilmington.  \citep[8-13-86]{ws}\par


\chapter{The Raft Runners}

\textit{\textsc{When} the Lord was making the world He had a lot of
  river left over, so he doubled it up and made it into the South
  (Black) River.  Thus, said the old folks, is how this stream came to
  be so crooked.}  \citep{shawg}\par

South River also is narrow.\par

Yet the small white farmers of Bladen, Sampson, and Cumberland
counties floated down their logs and barrels of naval stores to Black
River where they were formed into rafts and carried on to market in
Wilmington.\par

Both these and other Sampson County farmers who came down Six Runs and
the Coharies to the Black River were a hard working and hardy set of
people.  They came in winter and early spring when the streams were in
freshet and when the weather was cold and nasty.\par

To the people of the lower river they were strangers from a seemingly
vast and mysterious country.\par

Although tough and rugged they were strangers with some pleasant tune
on their lips except when fighting the wiles of the river, the
swirling currents of sharp bends, and the eddies braking and twisting
their train of logs.\par

These were strangers who were ever ready to hail a ``Hello there!'' to
any person who would return their greeting.\par

Some of the raft runners, several days from their upriver homes, seemd
a bit lonely.  While sitting beside their campfire or riding the river
current they dismissed a part of their solitude and at the same time
awakened lonely countrysides with hooting, hollering, singing, and
whistling.  The ddep river swamps amplified and echoed these sounds,
and at times they could be heard for many miles.  So those dwellers
beside the main dusty road, which followed near by the river, and
others afar might hear musical strains that belonged to earlier times
when raft runners of earlier generations passed this same way.\par

They rode rafts of many kinds.  From colonial times on past the Civil
War they carried to market choice timbers for masts and spars for
sailing ships, squared logs, or ton timber, for sawmills, and naval
stores, like turpentine, rosin, and tar in trains of barrels.  After
ascendance of the railroads cross ties were stacked, clamped, and
formed into sections like log bateaux to form rafts; these were
floated to the rail heads in Wilmington.  Log rafts upon large waters
might contain as many as five hundred logs with bateau after bateau
following in tandem.\par

\section[Stream Clearance Scheme]{A Stream Clearance Scheme}

By the opening of the twentieth century, rafting on the Black River
had gone on so long that it had become an institution.  Upon her
tributaries a hundred years earlier, people were questing after an
unobstructed water avenue to market.  Announcement of a lottery was
made in the April 5, 1811, \textit{Cape Fear Recorder}; and it stated
that the ``scheme was for the purpose of clearing and making navigable
for rafts and boats the main branch of the Six Runs in Sampson County,
from Kerby's Landing to Fryer's Bridge, so called.''
\citep[4-5-11]{cfr}.  Soon after 1900 and following heavy rains a
Harrells farmer floated cross ties down nearby Bulltail Swamp to Black
River and on to market \citep[]{melvinew}.  Narrow rafts, six or eight
logs wide, could be brought down South River, a long and crooked Black
River tributary, from the vicinity of Dunn in Harnett County
\citep[]{johnsonlg}.\par

Large Numbers of rafts came down the Black River and other streams of
North Carolina's southeast.  The December 17, 1852 Wilmington
\textit{Journal} stated, ``about 100 rafts of timber changed hands''
during the previous week.  Then on December 24 following it explained:
``The water courses have been high for some time, and produce such as
timber, Turpentine, Lumber \&c, usually dependent upon the state of
the rivers, has been coming in abundantly lately.''\par

\section[An Adventure]{An Adeventure of an Eight-Year-Old Boy}

Sometimes the farmers brought along their young sons for the adventure
and to teach them the ways of the river.  A vivid account of raft
running during winter has been given by Campbell Johnson, a Bladen
County boy who at the age of eight or ten years made a trip by raft to
Wilmington, a city of wonders for country children.\par

The loud cry of a steamboat whistle came thundering down the river,
from around a far-away bend of the the yellow water Cape Fear River.
It cried again and again warning any craft it might be meeting.  It
was a welcome and a thrilling call for the boy; for this was the first
familiar thing he had experienced on the new and strange lower river.
He had been hearing the same steam whistle as far back as he could
remember, for it frequently broke the monotony of the quiet
countryside near his home in Rowan community, Bladen County, about
forty miles upstream as the boat plied the Black River, a Cape Fear
tributary.\par

``The \textit{Liz}!'' the boy cried as he bounded from beneath an
inverted log boat sheltering him from the cold drizzling rain.\par

He was on a log raft basking before a roaring log fire.  His father
and two helpers had tied up their raft at Long Reach on the Cape Fear
to await the change of the tide which would give them a free ride
several miles more to the sawmill in Wilmington.  The \textit{Lisbon},
which some Black River folks called \textit{Old Liz}, was Captain
D.~J.\ Black's\index{Black, Captain D.~J.} river steamer.  She was
coming down from Clear Run, twelve miles below the town of Lisbon and
the confluence of Six Runs and Great Coharie.  Captain Black lived at
Point Caswell, and he and \textit{Old Liz} paddled up and down both
the Black and Cape Fear rivers picking up and delivering freight and
carrying a few passengers.\par

Soon after the \textit{Lisbon's} whistle blasts the thumping and
pounding of her steam engine floated down the river, and she came
around the far-away bend.  The boy fastened his eyes upriver like
folks waiting at river landings upon the approach of a steamer.  It
was practically no time before the \textit{Lisbon} drew near.  White
spray frothed about her low bow which seemed to plow a furrow in the
river like a farmer in his field.  The boat slowed, and as her
whistled signaled a greeting ``boop, boop, boop'' the captain waved
from his pilot house to the men and the boy on the raft.\par

\section{Hard Life on the River}

Although the hardy men along the rivers refused to let foul weather
hinder them in their raft running, they often talked about the
discomfort and hardships they suffered.\par

Except for a few rude shelters they were exposed to the weather.  Soon
after 1900, for example, Frank Sessoms, a storekeeper at Long View on
the Black River, provided some comfort for the raft men coming down
from Sampson and Cumberland Counties.  He had an old shanty for them
to camp out in when the weather was bad.  \citep[]{rawlsgw}\par

But, the Bladen County rafters had no shelter once they left
Sparkleberry near Mill Creek on the Black River before daybreak the
morning of the previous day.  At Sparkleberry where the river widens
the narrow long rafts which had run the \textit{Narrows} were widened.
The big raft now was actually a train of bateaux, a bateau being a
section of a dozen or more logs fastened side by side at each end with
batten poles.  In 1900 hickory pegs were driven into holes bored
through the batten poles into the logs.  A few years later sixty penny
size spikes replaced the wooden pegs.\par

Already, the rain had begun to come down in a cold drizzle.  The men
prepared for a hard trip.  They placed boards side by side across the
logs of the front bateau.  They covered the boards with a layer of
dirt, and this enabled them to build a fire.  They piled nearby a tall
stack of fire wood cut from small trees.\par

Before the ropes holding the raft fore and aft were untied and the
raft set upon its one to two mile-an-hour drift the fire was roaring.
But the rain grew heavier, and the boy stood near the fire as moisture
steamed from his wet clothes.  One of the three men steered the course
of the raft with a long home-made oar which extended ahead of the
front bateau.  A second man worked from section to section on the rear
bateaux making certain they ran clear of obstructions.  The third man
took his turn before the fire and was ready to help when and where
needed.\par

The wooden \textit{cracker box}\footnote{Crackers carried by country
  storekeepers came packed in wooden boxes.  The empty boxes found
  many uses; one was for packing food for picnics and as in this
  instance a trip of several days.}, filled with food for a trip of
three or four days, stood out in the rain.  The women folks back home
had packed it with home-made cookies known as \textit{tea cakes}, fat
crusty flour biscuits and boiled meats.  In addition, there were meal
for fresh corn bread, eggs, lard, and a pan for frying, fishing lines
and fat meat to entice catfish during hours when the raft was tied up
against the inflowing tide.\par

The raft reached Racoon Island near the mouth of the Black River the
two tides it rode the first day.  Here it was tied up for the night.
The men seemed tired from a long day's work.  After eating supper they
``wrapped quilts around their heads'' and lay around the open fire.
They went to sleep with apparent ease even though the rain streamed
down cold.\par

By the following afternoon the raft had ridden the tide to Long Reach
where it was tied up to await the next outflowing tide.  The bateaux
lay end to end beside the tall grass of the marshlands where only a
few years earlier there had been rich rice fields.  Sand banks rose up
high on the opposite side of the river; they were known as
\textit{Mount Misery}.  Except for scattered scrub oaks and pines the
place seemed desert-like.  Perhaps its massive desolation made it one
of the more impressive landmarks witnessed by those who traveled the
lower Cape Fear River.\par

Beyond the march grass where the land rose to higher ground, smoke
lazed from the mud and stick chimney of a small cabin and drifted in a
foggy layer towards the ground.  Openings between the logs in the
cabin's walls had been daubed with clay.  The cabin looked old and
time-worn like the old black man and his woman who made it their home;
but the cabin was warm inside, and the boy enjoyed the hospitality of
the plain hearth.\par

The old man tilled a small plot of farm land which rose above the
marshes; and his chickens and hogs enjoyed open range.\par

His hogs were of unusual appearance to the boy who was accustomed to
the wild razorbacks of the swamps and great woods and the red and
black spotted ones back home in pens and pastures.  These hogs were
Blue guineas, the old man said.  The boy begged his father to buy him
a small gilt.  The father thought and then said, ``Maybe we can get
her home'' and paid the two dollars the old man asked.\par

Later, at Wilmington, Captain Black\index{Black, Captain D.~J.} was
told of the purchase, and he said he could take the hog on his next
trip upriver.  Such cargo was not unusual for Black River steamboats.
Their masters were accustomed to hauling almost anything the farmers
wanted to ship, including chickens, eggs, hogs, sheep, and cattle.\par

Early the next morning the three raft runners and the boy caught the
incoming tide and sent their twenty-foot-long log boat skimming the
solitary waters for home.  Early that night they arrived at David's
Landing, one mile along the public road from home, having paddled
forty-eight miles.\par

The following afternoon the \textit{Lisbon} arrived with \textit{Liz}.
The boy had named the gilt in honor of
Captain Black's\index{Black, Captain D.~J.} boat.  \citep{johnsonwc}\par

\section{The Prized Piano}

Sometimes Sampson County farmers above the head of navigation for
Black River steamboats would carry their farm produce to market on
their squared timber or log rafts.  The produce, chiefly corn, cotton,
hay, and bacon was piled atop a decking of loose boards, which too
were sold with the raft.\par

This had been the custom of Haywood J.\ Peterson, a Sampson County
farmer who lived near Clinton.  At the age of sixty-two in 1891
Peterson decided to move to Point Caswell on the lower Black River and
go into semi-retirement.  He sold his farm, loaded his household
furniture on a raft, and this rode a freshet down the Great Coharie
and the Black River to its destination.  But Peterson's daughter,
Hyancith, who was to become the wife of Captain
D.~J.\ Black\index{Black, Captain D.~J.}, would not trust her prized
piano on the raft.  It was too dangerous, she insisted; so the piano
was carried overland by horse and wagon.  It was delivered safely,
although a bit shaken by the rough roads, to the Point Caswell Hotel
where the family would live until a dwelling could be built.\par

A lively freshet carried the raft down swiftly.  The water was so high
that in places it streamed over the river banks into the bordering
swamps.  The raft runners fought to avoid being drawn into the forests
and demolished.  At one point they lost control and the current
carried their raft beneath an overhanging limb which swept a part of
the furnishings into the river.  Those which floated were picked up
and saved; but the sinkables, such as the family's treasured
silverware, were lost.\par

Hyancith treasured her piano until her death in 1976.\\
\citep{shumpertsb}.\par

\section{The Snow Train}

Once while recounting the rigors of raft running, one tough lumberjack
told of a trip down the Waccamaw River.  He was drifting a raft of
about 300 logs by himself, and a blinding snowstorm forced him to tie
up before nightfall.\par

He went ashore, gathered a quantity of fire wood, constructed a
lean-to shelter and made a bed of pin straw and Spanish moss.\par

He built a large fire in front of the shelter, ate his supper and
slept warmly while the wind howled and the snow came down.  By morning
the snow was about two feet deep, and his shelter was much like an
igloo.  His raft, decked in white, contrasted sharply with the black
water in the river.\par

The storm had mostly blown over, and he resumed his journey with the
current.  Two days later he arrived at the mill where no snow had
fallen.  He became the center of attention as a crowd of people lined
the river bank to admire his long procession of snow clad logs.
\citep{stevensjg}.\par

\makeatletter\@openrightfalse%

\part{II}

\begin{figure}[htb]
\centering
  \begin{turn}{90}
      \begin{minipage}{5.8in}
        \centering
        \includegraphics[scale=.18]{ballous_pictorial_24_feb_1855_pp_121}
        \caption{Market St Wilmington}
      \end{minipage}
    \end{turn}
\end{figure}

\chapter{The First Steamboats}
\@openrighttrue\makeatother

\section{The Monopolies}

\textsc{The} year 1818 marked the beginning of steamboat navigation in
eastern North Carolina when two paddle wheel steamers were put into
service on the Cape Fear River and a third one on the Neuse.\par

The Cape Fear steamers where the \steamer{Henrietta}, which was built
at Fayetteville by James Seawell\index{Seawell, James} to run between
that point and Wilmington, and the \steamer{Prometheus}, which was
built at Beaufort by Captain Otway Burns\index{Burns, Captain Otway},
famed as master of the privateer
\textit{Snap~Dragon}\index{Snap~Dragon,~privateer}, for a stock
company, which may have been the Cape Fear Navigation Company.  The
\steamer{Prometheus} was put in service on the lower river between
Wilmington and Smithville, now Southport.\par

The Neuse steamer was the \steamer{Norfolk}, built out of state and
purchased by the New Bern Steamboat Company.  Thomas H. Sloan in his
\textit{Inland Steam Navigation in North Carolina 1818-1900}, citing % chktex 8
the New Bern \textit{Centinel} of April 11, 1818, states that the
\steamer{Norfolk} arrived in New Bern April 10 and thus became the
first steamer to go into service in the state.  Twenty days later, on
April 30, according to the May 9 Centinel, the \steamer{Henrietta} was
launched on the Cape Fear; and the July 23
\textit{National Intelligencer} of Washington D.~C., tells of her
first trip to Wilmington.  The June 13 \textit{Intelligencer} had
stated that the \steamer{Prometheus} had arrived at Wilmington and
that within a few days she would probably commence running from that
place to Smithville.\par

These reports seem to establish the \steamer{Norfolk} as the first
steamboat service in North Carolina and the \steamer{Henrietta} as the
first steamship built in this state.\par

However, unsuccessful steps to introduce steamboat service had been
taken several years earlier.  Robert Fulton\index{Fulton,~Robert} had
made his first successful voyage on the Hudson River with his
steamboat \steamer{North~River}, otherwise known as
\textit{Claremont}\index{Claremont|see {North~River}}, eleven years
before the \steamer{Henrietta} and the \steamer{Prometheus} began
service on the Cape Fear.\par

A proposal to bring steamboat service to North Carolina was made by
John Stevens, of Hoboken, New Jersey, in 1812 when he presented a
memorial to the State Legislature.\par

In response an act was passed ``for the encouragement of Steam Boats
on the waters of this State.''  By this act, Stevens and his
associates were given a twenty year monopoly for steamboat operation.
But, the act provided that Stevens must have a boat in operation
within two years for the monopoly to continue on force.
\citep[12]{sloanth}.\par

Stevens' monopoly was met with several protests, one being made by
Fulton.  Even though his monopoly was declared valid by the
legislature, Stevens failed to establish steamboat service as
provided.  It lapsed in 1814, and the legislature refused to grant his
request for renewal even though he explained that he had built the
steamer \textit{Phoenix} for this service and the British blockade of
the coast during the War of 1812 had forced her to remain in the
harbor at Hoboken.  \citep[1-3]{turlingtonsw}.\par

The protests and ensuing debates in the legislature for the improved
transportation that steamboats would provide aroused widespread
interest in steam navigation.  One of the main obstacles was the poor
condition of the rivers.  The presence of logs and snags made
navigation extremely difficult and hazardous.  However, steps were
taken to remedy the difficulties.  The Cape Fear Navigation Company
had its original stock supplemented with \$50,000 in private
subscriptions.  Work began on the river before august 1816, and
Archibald Murphey stated in 1819 that, ``the river is navigable for
Steam-Boats to Fayetteville, and may be probably rendered navigable
for small Steam-Boats to Haywood.''  \citep[7, 11]{turlingtonsw}.\par

By 1838 the Cape Fear Navigation Company had removed 1159 trees, logs,
and stumps from the river.  Wing dams had been built, and navigation
on the Cape Fear was considered ``better, safer, less liable to
interruption from low water, and by far cheaper than any river south
of the Potomac, of the same size and extent.''  By this time the
steamboat \steamer{Henrietta} had run the river for twenty years
without being snagged or meeting with any serious accident.
\citep[12, 13]{turlingtonsw}.\par

After turning down a petition in 1817 by Joseph
Seawell\index{Seawell, Joseph} of Fayetteville, the General Assembly
of the following year granted him a monopoly for steam navigation
between Wilmington and Fayetteville.  His petition stated that he had
a boat under construction.\par

The bill gave Seawell\index{Seawell, Joseph} and his associates ``the
sole and exclusive right and privilege of using, employing, and
navigating all and every species or kinds of boats propelled by the
force of steam on the waters of the Cape Fear river between the towns
of Fayetteville and Wilmington for and during the full term of seven
years from and after the passage of this act.''
\citep[22]{turlingtonsw}\citep[13]{sloanth}.\par

Launching of the \steamer{Henrietta} at Fayetteville and arrival of
the \steamer{Prometheus} at Wilmington were met with great
enthusiasm.\par

\section{The \textit{Prometheus}}

The \steamer{Prometheus} was brought down from Beaufort by Captain
Thomas N.~Gautier\index{Gauitier, Captain Thomas N.}, a worthy old sea
captain.  Colonel J.~G.~Burr told the story of her arrival to James
Sprunt who gives it in his
\textit{The Chronicles of the Cape Fear River}.
\citep[139]{spruntj1916}.\par

``Expectations were on tiptoe after the departure of the captain; a
feverish excitement existed in the community, which daily increased,
as nothing was heard from him for a time, owing to the irregularity of
the mails; but early one morning this anxiety broke into the wildest
enthusiasm when it was announced that the \steamer{Prometheus} was in
the river and had turned the Dram Tree.  Bells were rung, cannon
fired, and the entire population, without regard to age, sex, or
color, thronged the wharves to welcome her arrival.  The tide was at
the ebb, and the struggle between the advancing steamer and the fierce
current was a desperate one; for she panted fearfully, as though
wind-blown and exhausted.  She could be seen in the distance,
enveloped in smoke, and the scream of her high-pressure engine
reverberated through the woods, while she slowly but surely crept
along.  As she neared Market Dock, \ldots the captain called through
his speaking trumpet to the engineer below: `give it to her, Snyder';
and while Snyder gave her all the steam she could bear, the laboring
\steamer{Prometheus} snorted by, amid the cheers of the excited
multitude.''  \citep[139]{spruntj1916}.\par

The \steamer{Prometheus} was in operation in 1819 when President James
Monroe and his party were hospitably entertained in Wilmington.  On
his departure the President was taken to Smithville on the new
steamboat.  \citep[4-27-1819]{ni}\par.

But the low-powered \steamer{Prometheus} seems to have been less than a
financial success; she was abandoned in 1825, only seven years after
her launching.  \citep[19]{sloanth}.\par

\section{The \textit{Henrietta}}

But the story of the \steamer{Henrietta} was quite a different one
from that of the \steamer{Prometheus}.  As a trail blazer for other
pioneer steamboats on the Cape Fear River her long life and financial
success was remarkable.\par

She was built by James Seawell\index{Seawell, James} in 1817-18, on % chktex 8
his plantation on the east side of the Cape Fear River three miles
above Campbellton, now Fayetteville.  Seawell\index{Seawell, James}
named the boat for his daughter Henrietta.  A large concourse of
people attended her launching on April 30, 1818.  Then she was put
under command of Captain Charles Taws\index{Taws, Captain Charles}.
\citep[429]{oatesja}.\par

Many flat, or pole, boats had been constructed in the Campbellton
area, and the \steamer{Henrietta}'s sturdy hull must have profited
from the skill of local ship carpenters.  \citep[428]{oatesja}.\par

She was a side-wheeler which drew seven and one-half feet of water,
\citep[342]{oatesja}; a strange looking craft and something of an ugly
duckling when compared to later steamers.  \citep[7-24-1891]{ws}.\par

In 1891 old Bob Orrell\index{Orrell, Captain Bob}, who had retired
from the river and for many years on the decks of the
\steamer{Henrietta}, recalled, ``She had no upper deck at first, and
her cabin was set down in the hold.''  She looked much like the dingy,
gunwale gnawed, splintery steam flats which frequented the Cape Fear
River and her tributaries after the Civil War.
\citep[7-24-1891]{ws}.\par

The \steamer{Henrietta} made her first trip from Fayetteville to
Wilmington in early July 1818.  Under the Fayetteville dateline of
July 16 the \textit{National~Intelligencer} of July 23 announced that
``the steam boat \steamer{Henrietta} left this place on Saturday last,
at the rate of about 8 miles per hour, for Wilmington, where she
arrived on Sunday without sustaining the least damage---not
withstanding she met with numerous impediments from the present
lowness of the water.''  \citep[7-23-1818]{ni}.\par

Then under the Wilmington dateline of July 18, the
\textit{National Intelligencer} of July 28 reported, ``The Steam Boat
\steamer{Henrietta}, of Fayetteville, arrived at this place on Sunday
last.''  The 115 mile trip down stream had taken six days; but her
movements on the lower river, between Wilmington and Smithville, were
``very rapid.''  It was noted that ``much credit is due to those who
furnished so excellent a vessel for the accomodation of the citizens
of Fayetteville, Wilmington, and Smithville.''  \citep[7-28-1818]{ni}.
Although the \steamer{Prometheus} is not mentioned, the liberal praise
of the \steamer{Henrietta} leads one to believe that she was
recognized as a superior steamboat.\par

Nonetheless, from the outset the \steamer{Henrietta} was beset with
mechanical problems.  She was geared ``to work with cog wheels like a
mill'' and when she reached a dozen or more sharp points in the river
she had to be \textit{dropped} around with a line.  That was almost as
bad as trying to sail a schooner or some other large sailing ship up
the narrow and crooked river.  Often some of the crew had to put out
in a small boat and secure a line to a tree or some other well
anchored object so that those on her bow could hand pull or windlass
her around while others steered her from shoals with long poles.\par

Although the river had been cleared of obstructions by 1819, the
\steamer{Henrietta}'s annoying mechanical problems persisted; and this
brought Captain Benjamin Rush\index{Rush, Captain Benjamin} down from
Philadelphia in 1820.  At first Rushed changed her gearing to a chain
motion, and when this did not remedy the problem he changed it to a
connecting rod and crank motion.  This enabled her to round the sharp
points with no great difficulty.  Rush continued on as the boat's
master.  In appreciation of the engineer's achievement and service one
of the \steamer{Henrietta}'s large tow boats
was given his name, the \barge{Ben Rush}.\par

With service made more dependable the \steamer{Henrietta}'s passenger
and freight service increased.  So, she was given an upper deck and
better accommodations for more passengers.  \citep[7-24-1891]{ws}.\par

Rush continued with the \steamer{Henrietta} most of her forty years on
the Cape Fear River.  Duncan MacRae of Fayetteville thought him to be
a model of the versatile old river captain.  He described him as ``a
gentleman, affable, attentive, and popular with all.''
\citep[91]{oatesja}.  A notice in 1838 listed Rush as part owner of
the steamer.\par

Bob Orrell\index{Orrell, Captain Bob} said that the
\steamer{Henrietta} was kept modernized and profitable.  She was given
powerful machinery to bring down large cargoes of naval stores,
lumber, and other freight from upriver, chiefly Fayetteville but
frequently Whitehall and Elizabethtown, at good speed.  She might be
seen loaded to her gunwales panting along towing two or three large
flat boats.\par

When Bob Orrell\index{Orrell, Captain Bob} was managing the Henrietta
Steamboat Line\index{Henrietta~Steamboat~Line} he voiced the opinion
that the \steamer{Henrietta}, kept in prime condition in her old age,
was the fastest boat on the river.  She had run the 115 miles from
Fayetteville to Wilmington in less than ten hours.\footnote{The
  Raleigh Register of March 14, 1837, reported that the steamer
  \steamer{Wilmington} had made it in seven and one half hours.}
This, some of the river steamer captains said, was not the case.  So
Orrell\index{Orrell, Captain Bob} and Captain
Lammerman\index{Lammerman,~Captain} agreed to a race between the
side-wheel \steamer{Henrietta} and the newer twin-boiler, stern-wheel
\steamer{Chatham}.  Word of the contemplated race created considerable
interest among the river men and Fayetteville townsmen.\par

The race would be from Fayetteville to Wilmington.  Passengers made
their choice of boats, and the \steamer{Chatham} was sent on her way
fifteen minutes before the \steamer{Henrietta} left her Fayetteville
berth.  A scene similar to this may have followed:\par

A few miles down the Cape Fear, the \steamer{Henrietta}'s passengers
saw the \steamer{Chatham} come into view.  Now, streams of smoke
poured from the dual stacks of both steamers; their paddle wheels
clawed at the river.  The \steamer{Henrietta} slowly closed the
\steamer{Chatham}'s lead; then, on a straight-of-way with engines
throbbing, steam hissing, passengers and whistles screaming, the
twin-wheeled \steamer{Henrietta} snail-paced past the
\steamer{Chatham} and left her churning the former's wake.\par

Orrell\index{Orrell, Captain Bob} said the \steamer{Henrietta} had
left Fayetteville ``at exactly 6 o'clock a.~m., and my passengers were
uptown in Wilmington at 4:30 in the afternoon.''  She had made only
made two stops for wood and had beaten the \steamer{Chatham} by two
hours.\par

But Orrell\index{Orrell, Captain Bob} was charitable; no other boat on
the river at that time could beat the \steamer{Chatham}.  She was one
of the few keel-bottomed boats, had high horsepower engines, and was
very fast when loaded.  In time her great power led her to
destruction.  She blew her boiler in 1853.\par

The \steamer{Henrietta} was the first steamer to go up the Cape Fear
River as far as Averasboro in Harnett County.  This was long before
regular steamer service commenced in 1856.  Her smokestack was removed
so she could pass under the Clarendon bridge.  She took along three
flat boats and brought back 3,400 barrels of rosin.\par

On one occasion, Orrell\index{Orrell, Captain Bob} recalled, the
\steamer{Henrietta} with the help of the big tow boat \barge{Ben
  Rush}, to two Fayetteville merchants 4,500 bushels of rock alum salt
and 95 hogsheads of molasses.  These merchants had established a large
trade with people of the North Carolina piedmont.\par

Orrell\index{Orrell, Captain Bob} estimated that when the
\steamer{Henrietta} retired from regular service she had run at least
a million and a half miles on the Cape Fear River and she had earned
for her owners no less than a million and a half dollars.\par

At times during the \steamer{Henrietta}'s long life, in addition to
Captains Taws\index{Taws, Captain Charles} and
Rush\index{Rush, Captain Benjamin}, she was commanded by Captains
Orrell\index{Orrell, Captain Bob},
Wilkerson\index{Wilkerson, Captain},
Hartmann\index{Hartmann, Captain}, and
Rothwell\index{Rothwell, Captain}.  \citep[7-24-1891]{ws}
\citep[Note on p.~139]{spruntj1916} \citep[179, 192-193, 342-342]{oatesja}.\par

Soon after the \steamer{Henrietta} was put into service, James
Seawell\index{Seawell, James} built a second steamboat at Fayetteville
for a stock company which had for its major stockholder Duncan
Thompson\index{Thompson, Duncan}.  She was the first of three steamers
to bear the name \steamer{North Carolina}.  By early 1823, the
\steamer{North Carolina} had joined the \steamer{Henrietta} in
rendering freight and passenger service on the Cape Fear River between
Wilmington and Fayetteville.  \citep[4-27-1823]{hr}.  Captain Charles
Taws\index{Taws, Captain Charles}, first master of the
\steamer{Henrietta} was given command.  Three flats added to the
freighting capacity of the two steamers.  \citep[10-23-1823]{co}.\par

There is no mention of the \steamer{North Carolina} having major
mechanical difficulties.  Captain Ben
Rush's\index{Rush, Captain Benjamin} troubleshooting on the \steamer{Henrietta}
must have profited in the construction of the new boat.\par

In 1825, one finds the \steamer{North Carolina} running excursions.
The Fayetteville Observer reported that large numbers of persons
boarded her, and she took them several miles to the vicinity of where
a camp meeting was being held.  \citep[394]{johnsongg}.\par

As steam navigation became established on the Cape Fear, merchants of
the interior were urged to patronize the Wilmington harbor.
Advantages cited were West Indies produce was cheaper at Wilmington
than any port in the Atlantic States; forest products could be used in
exchange for salt and other products; and it was stated that North
Carolina's shipping tonnage was 23,760 compared with 16,301 for
Virginia.  \citep[26]{turlingtonsw}.\par

The \textit{Hillsborough Recorder} pointed out that ``There are two
steam boats, besides a number of poll boats and flats, constantly
running on the Cape Fear, between Fayetteville and Wilmington, and the
prices generally run about the same, with the addition of freight,
which is only about 20 cts.\ per 100 lbs.  For some years past, the
Fayetteville market has been generally well supplied with groceries;
there are also several houses that import dry goods, and as the demand
increases, the supplies and the varieties will increase in the same
ratio.  In several instances have goods been received here
(Fayetteville) on the seventh day, after they were shipped from New
York.''  \citep[4-2-1823]{hr}.\par

\chapter{Ascendancy of Steam Navigation}
\chaptermark{Steam's Ascendancy}

\textsc{When} James Seawell's\index{Seawell, James} seven-year
monopoly ended 1 January 1825, other interests were prepared to share
in the lucrative steamboat business.\par

Seawell and his associates prepared for the competition.  The
Henrietta Steamboat Company\index{Henrietta~Steamboat~Company} was
incorporated in the town of Fayetteville by Horatio
G. Nelson\index{Nelson,~Horatio~G.} and Benjamin
Rush\index{Rush, Captain Benjamin} of Fayetteville, and Doyle
O'Hanlon\index{O'Hanlon,~Doyle}, the ship builder of Wilmington.  The
company was authorized to build one or more additional steamboats and
the necessary tow boats.  It was also authorized to hold land along
the water fronts at Fayetteville and Wilmington and along the river
that wharves and warehouses might be built.
\citep[27]{turlingtonsw}(N.~C.~Laws, 1827, 37, 38).\par

The Cotton Plant Steamboat
Company\index{Cotton~Plant~Steamboat~Company} was chartered in 1827
with the same privileges for their steamer, the
\steamer{Cotton~Plant}, \citep[27]{turlingtonsw}, which had been built
in 1826 preparatory to commencing service the following year.  The
steamer was a side-wheeler of the old walking beam type.\par

By now cotton production, stimulated by the general utilization of the
cotton gin, accounted for freight second only to naval stores, and it
seems that the boat's owners found it proper to name her for the plant
which promised to be so important to her.\par

Before her maiden trip to Wilmington ``a large party of (Fayetteville)
ladies and gentlemen'' came aboard the ``elegant'' steamer and were
treated to an excursion down the Cape Fear River ``while music,
dancing, and other diversions, and a profusion of good cheer,
enlivened the scene.''  \citep[151]{johnsongg}.  For days after her
arrival in Wilmington, crowds of people overflowed her decks.\par

The steamer enjoyed twenty or more years active service.  As late as
1847, she was on a regular schedule between Wilmington and
Fayetteville under command of Captains John
Banks\index{Banks, Captain John} and Thomas M. Green\index{Green, Thomas M.}.
\citep[Marine Intelligence]{wj}.  Years later she sank just below the
Railroad wharf.  Her hull was abandoned but her machinery was salvaged.
\citep[clip, n.d.]{fo}.\par

Then came the ill-fated \steamer{John Walker}, built by built by the
O'Hanlon Ship Yard\index{O'Hanlon Ship Yard} and named for John
Walker, father of Dr.~John Walker and others of Wilmington.  This
steamer was first built at Wilmington.  In 1831 her boiler exploded
below Wilmington, killing Captain Dickson\index{Dickson, Captain} and
her engineer.  \citep[11-23-1831]{cfr}\citep[2-10-1882]{ws}.\par



\begin{thebibliography}{99}

\section*{Books, Pamphlets, Theses, Papers}

\bibitem[\protect\citeauthoryear{Brickel,~J.}{1968}]{brickellj}
  Brickell, John.  \emph{The Natural History of North Carolina}, 1968
  reprint of 1737, Johnson Publishing Co., Murfreesboro, N.~C.

\bibitem[\protect\citeauthoryear{Clark,~W.}{1861-1865}]{clarkw}
  Clark, Waler.  \emph{North Carolina Regiments, 1861-1865}, % chktex 8
  5 volumes, Volume IV. % chktex 13

\bibitem[\protect\citeauthoryear{Crittenden,~C.~C.}{April 1931}]{crittenden}
  Crittenden, Charles Christopher.  ``Inland Navigation in North Carolina, 1763-1789'', % chktex 8
  \emph{North Carolina Historical Revue}, VIII (April 1931), 145-154. % chktex 8

\bibitem[\protect\citeauthoryear{Evans,~W.~McKee}{1966}]{evansw}
  Evans, W. McKee.  \emph{Ballots and Fence Rails, Reconstruction on the Lower Cape Fear}.
  Chapel Hill, University of North Carolina Press, 1966.

\bibitem[\protect\citeauthoryear{Fries,~A.~L. and others}{}]{friesal}
  Fries, A.~L.~and others.  \emph{Records of the Moravians in N.~C.},
  12 volumes, Vol. I.

\bibitem[\protect\citeauthoryear{Johnson,~G.~G.}{1937}]{johnsongg}
  Johnson, Guion Griffis.  \emph{Ante-Bellum North Carolina, A social History}.
  Chapel Hill, University of North Carolina Press, 1937.

\bibitem[\protect\citeauthoryear{Lee,~L}{1937}]{leel}
  Lee, Lawrence.  \emph{The Lower Cape Fear in Colonial Days}.
  Chapel Hill, University of North Carolina Press, 1937.

\bibitem[\protect\citeauthoryear{Lefler,~H.~T.}{1956}]{leflerht}
  Lefler, Hugh Talmadge.  \emph{History of North Carolina}, 4 volumes.  New
  York, Lewis Publishing Company, 1956.

\bibitem[\protect\citeauthoryear{Oates,~J.~A.}{1950}]{oatesja}
  Oates, John A.  \emph{The Story of Fayetteville and the Upper Cape Fear}.
  Charlotte, The Dowd press, 1950.

\bibitem[\protect\citeauthoryear{Olmstead,~F.~L.}{1904}]{olmsteadfl}
  Olmstead, Frederic Law.  \emph{A Journey in the Seaboard Slave States, in the
    Years 1853-1854, With Remarks on Their Economy}, 2 volumes.  New York and % chktex 8
  London, The Knickerbocker Press, 1904.

\bibitem[\protect\citeauthoryear{Powell,~W.~S.}{1968}]{powellws}
  Powell, William S.  \emph{The North Carolina Gazeteer}.  Chapel Hill,
  University of North Carolina Press, 1968.

\bibitem[\protect\citeauthoryear{Reilly,~J.~S.}{1884}]{reillyjs}
  Reilly, J.~S.  \emph{Wilmington---Past, Present, and Future}.  Wilmington,
  1884.

\bibitem[\protect\citeauthoryear{Roberson,~E.~E.}{}]{robersonee}
  Roberson, Elizabeth Ellis.  \emph{Diary}.  Elizabethtown, Bladen County Historical Association.

\bibitem[\protect\citeauthoryear{Schaw,~J.}{1939}]{schawj}
  Schaw, Janet.  \emph{Journal of a Lady of Quality}.  New Haven, Yale
  University Press, 1939.  Spartanburg, S.~C.~reprint.

\bibitem[\protect\citeauthoryear{Sharpe,~B.}{1954, 1958, 1961, 1965}]{sharpeb}
  Sharpe, Bill.  \emph{A New Geography of North Carolina}, 4 volumes.  Raleigh,
  Edwards \& Broughton Co., 1954, 1958, 1961, 1965.

\bibitem[\protect\citeauthoryear{Sloan, T.~H.}{1971}]{sloanth}
  Sloan, Thomas H.  \emph{Inland Steam Navigation in North Carolina 1818-1900}. % chktex 8
  Unpublished master's thesis, East Carolina University, Greenville, N.~C., 1971.

\bibitem[\protect\citeauthoryear{Sprunt, J.}{1916}]{spruntj1916}
  Sprunt, James.  \emph{Chronicles of the Cape Fear River 1660-1916}, Second % chktex 8
  Edition.  Raleigh, Edwards \& Broughton Printing Co., 1916.

\bibitem[\protect\citeauthoryear{Sprunt, J.}{1896}]{spruntj1896}
  Sprunt, James.  \emph{Tales and Traditions of the Lower Cape Fear 1661-1896}. % chktex 8
  Wilmington, LeGwin Brothers, 1896.

\bibitem[\protect\citeauthoryear{Turlington, S.~W}{1933}]{turlingtonsw}
  Turlington, Sarah Woodall.  \emph{Steam Navigation in North Carolina Prior
    to 1860}.  Unpublished master's thesis, University of North Carolina, Chapel
  Hill, 1933.

\bibitem[\protect\citeauthoryear{Whedbee, W.~L.}{}]{whedbeewl}
  Whedbee, W.~L.  \emph{Waterway Transportation in Eastern North Carolina}.
  Typescript, University of North Carolina Library.


\section*{Newspapers and Periodicals}

\bibitem[\protect\citeauthoryear{BJ}{}]{bj}
  BJ---\emph{Bladen Journal}, Elizabethtown, N.~C.

\bibitem[\protect\citeauthoryear{CC}{}]{cc}
  CC---\emph{Clinton Caucasian} Clinton

\bibitem[\protect\citeauthoryear{CFR}{}]{cfr}
  CFR---\emph{Cape Fear Recorder} Wilmington, N.~C.  Thomas Loring.

\bibitem[\protect\citeauthoryear{CO}{}]{co}
  CO---Not listed

\bibitem[\protect\citeauthoryear{FA}{}]{fa}
  FA---\emph{Fayetteville American}

\bibitem[\protect\citeauthoryear{FDC}{}]{fdc}
  FDC---\emph{Fayetteville Daily Courier}

\bibitem[\protect\citeauthoryear{FNCA}{}]{fnca}
  FNCA---\emph{Fayetteville North Carolina Argus}

\bibitem[\protect\citeauthoryear{FO}{}]{fo}
  FO---\emph{Fayetteville Observer}

\bibitem[\protect\citeauthoryear{FWM}{}]{fwm}
  FWM---\emph{Free Wilmington Magazine}

\bibitem[\protect\citeauthoryear{HR}{}]{hr}
  HR---\emph{Hillsboro Record}

\bibitem[\protect\citeauthoryear{MR}{}]{mr}
  MR---\emph{Magnolia Record}

\bibitem[\protect\citeauthoryear{NI}{}]{ni}
  NI---\emph{National Intelligencer}, Washington, D.~C.

\bibitem[\protect\citeauthoryear{NBCC}{}]{nbcc}
  NBCC---\emph{New Bern Carolina Centinel}

\bibitem[\protect\citeauthoryear{N&O}{}]{no}
  N\&O---\emph{News and Observer}, Raleigh, N.~C.

\bibitem[\protect\citeauthoryear{PC}{}]{pc}
  PC---\emph{Pender Chronicle}, Burgaw, N.~C.

\bibitem[\protect\citeauthoryear{RR}{}]{rr}
  RR---\emph{Raleigh Register and N.~C.~Gazette}

\bibitem[\protect\citeauthoryear{SA}{}]{sa}
  SA---\emph{Sampsonian}, Clinton, N.~C.

\bibitem[\protect\citeauthoryear{SL}{}]{sl}
  SL---\emph{Southport Leader}

\bibitem[\protect\citeauthoryear{WA}{}]{wa}
  WA---\emph{Wilmington Advertiser}

\bibitem[\protect\citeauthoryear{WFCR}{}]{wfcr}
  WFCR---\emph{Wilmington Cape Fear Recorder}

\bibitem[\protect\citeauthoryear{WDD}{}]{wdd}
  WDD---\emph{Wilmington Daily Dispatch}

\bibitem[\protect\citeauthoryear{WDH}{}]{wdh}
  WDH---\emph{Wilmington Daily Herald}

\bibitem[\protect\citeauthoryear{WDJ}{}]{wdj}
  WDJ---\emph{Wilmington Daily Journal}

\bibitem[\protect\citeauthoryear{WDP}{}]{wdp}
  WDP---\emph{Wilmington Daily Post}

\bibitem[\protect\citeauthoryear{WDS}{}]{wds}
  WDS---\emph{Wilmington Daily Star}

\bibitem[\protect\citeauthoryear{WHU}{}]{whu}
  WHU---\emph{Wilmington Herald of the Union}

\bibitem[\protect\citeauthoryear{WJ}{}]{wj}
  WJ---\emph{Wilmington Weekly Journal}

\bibitem[\protect\citeauthoryear{WN}{}]{wn}
  WN---\emph{Wilmington News}

\bibitem[\protect\citeauthoryear{WPP}{}]{wpp}
  WPP---\emph{Wilmington People's Press}

\bibitem[\protect\citeauthoryear{WC}{}]{wc}
  WC---\emph{Wilmington Weekly Commercial}

\bibitem[\protect\citeauthoryear{WM}{}]{wm}
  WM---\emph{Wilmington Weekly Messenger}

\bibitem[\protect\citeauthoryear{WS}{}]{ws}
  WS---\emph{Wilmington Weekly Star}

\bibitem[\protect\citeauthoryear{WSWM}{}]{wswm}
  WSWM---\emph{Wilmington Semi-Weekly Messenger}

\bibitem[\protect\citeauthoryear{WSMH}{}]{wsmh}
  WSMH---\emph{Wilmington Sunday Morning Herald}

\section*{Traditional Sources}

\bibitem[\protect\citeauthoryear{Arnold, W.}{}]{arnoldw}
  Arnold, Wayne, Burgaw, N.~C., Pender County.

\bibitem[\protect\citeauthoryear{Black, J.}{}]{blackj}
  Black, Jasper, Wilmington, N.~C., native of Pender County

\bibitem[\protect\citeauthoryear{Brown, I.~J.}{}]{brownij}
  Brown, Inez Johnson, Bladen County, N.~C.

\bibitem[\protect\citeauthoryear{Brown, V.}{}]{brownv}
  Brown, Vance, Bladen County, N.~C.

\bibitem[\protect\citeauthoryear{Campbell, W.~S.}{}]{campbellws}
  Campbell, Wanda S., Elizabethtown, N.~C., native of Bladen County, N.~C.

\bibitem[\protect\citeauthoryear{Cavalier, I.~L.}{}]{cavalieril}
  Cavalier, Irma Lewis, Long Creek, native of Pender County, N.~C.

\bibitem[\protect\citeauthoryear{Corbett, E.~J.}{}]{corbettej}
  Corbett, Evelyn Johnson, Jacksonville N.~C., Bladen County, N.~C., native.

\bibitem[\protect\citeauthoryear{Corbett, H.}{}]{corbetth}
  Corbett, Hill, Jacksonville, N.~C., native of Bladen County, N.~C.

\bibitem[\protect\citeauthoryear{Fisher, R.}{}]{fisherr}
  Fisher, Robert, Rocky Mount, N.~C., Native of Bladen County, N.~C.

\bibitem[\protect\citeauthoryear{Herring, N.~J.}{}]{herringnj}
  Herring, N.~J., Lewisburg, W.~Va., native of Sampson County, N.~C.

\bibitem[\protect\citeauthoryear{Highsmith, L.}{}]{highsmithl}
  Highsmith, Louis, Ingold, Route 4, Clinton, N.~C., native of Sampson
  County, N.~C.

\bibitem[\protect\citeauthoryear{Horrell, R.}{}]{horrellr}
  Horrell, Roland, Atkinson, N.~C., native of Pender County, N.~C.

\bibitem[\protect\citeauthoryear{Huffman, J.}{}]{huffmanj}
  Huffman, James, Delco, N.~C., native of Columbus County, N.~C.

\bibitem[\protect\citeauthoryear{Johnson, C.}{}]{johnsonc}
  Johnson, Clyde, Delco, N.~C., native of Bladen County, N.~C.

\bibitem[\protect\citeauthoryear{Johnson, L.~G.}{}]{johnsonlg}
  Johnson, Lee G., Watha, N.~C., native of Bladen County, N.~C.

\bibitem[\protect\citeauthoryear{Johnson, N.~K.}{}]{johnsonnk}
  Johnson, Nolon K., FRD, Ivanhoe, N.~C., native of Bladen County, N.~C.

\bibitem[\protect\citeauthoryear{Johnson, S.~H.}{}]{Johnsonsh}
  Johnson, Sadie Hill, Tampa, Florida, former resident of Bladen County, N.~C.

\bibitem[\protect\citeauthoryear{Johnson, S.}{}]{johnsons}
  Johnson, Sherwood, Rich Square, N.~C., native of Columbus County, N.~C.

\bibitem[\protect\citeauthoryear{Johnson, W.~C.}{}]{johnsonwc}
  Johnson, W.~Campbell, Tampa, Florida, native of Bladen County, N.~C.

\bibitem[\protect\citeauthoryear{Kelly, A.}{}]{kellya}
  Kelly, Aaron, RFD Kelly, N.~C., at Longview, native of Bladen County, N.~C.

\bibitem[\protect\citeauthoryear{Kelly, J.}{}]{kellyj}
  Kelly, John, RFD, Ivanhoe, N.~C., native of Bladen County, N.~C.

\bibitem[\protect\citeauthoryear{Kelly, M.}{}]{kellym}
  Kelly, Maurice, Atkinson, N.~C., native of Pender County, N.~C.

\bibitem[\protect\citeauthoryear{Melvin, E.~W.}{}]{melvinew}
  Melvin, E.~W.

\bibitem[\protect\citeauthoryear{Newkirk, A}{}]{newkirka}
  Newkirk, Armenius, Atkinson, N.~C., native of Pender County, N.~C.

\bibitem[\protect\citeauthoryear{Parker, J.}{}]{parkerj}
  Parker, Jim, in \textit{Sampsonian}, Clinton, N.~C.

\bibitem[\protect\citeauthoryear{Parramore, T.~C.}{}]{parramoretc}
  Parramore, Thomas C., Raleigh, N.~C., native of Hertford County, N.~C.

\bibitem[\protect\citeauthoryear{Pate, H.}{}]{pateh}
  Pate, Herb, in \textit{Pender Chronical}, Burgaw, N.~C.

\bibitem[\protect\citeauthoryear{Peterson, G.~J}{}]{petersongj}
  Peterson, Gertrude Johnson, Kelly, N.~C., native of Bladen County, N.~C.

\bibitem[\protect\citeauthoryear{Peterson, R.~F.}{}]{petersonrf}
  Peterson, R.~Fennel, native of Bladen County, N.~C.

\bibitem[\protect\citeauthoryear{Pridgen, S.}{}]{pridgens}
  Pridgen, Smith, Atkinson, N.~C., native of Bladen County, N.~C.

\bibitem[\protect\citeauthoryear{Rawls, G.~W.}{}]{rawlsgw}
  Rawls, George William, RFD, Currie, N.~C., native of Bladen County, N.~C.

\bibitem[\protect\citeauthoryear{Reaves, W.~E.}{}]{reaveswe}
  Reaves, William E., Wilmington, N.~C., New Hanover County

\bibitem[\protect\citeauthoryear{Robbins, L.}{}]{robbinsl}
  Robbins, Mrs.~Leo, Magnolia, Duplin County, N.~C.

\bibitem[\protect\citeauthoryear{Shaw, G.}{}]{shawg}
  Shaw, Graham, Atkinson, N.~C., native of Sampson County, N.~C.

\bibitem[\protect\citeauthoryear{Smith, T.}{}]{smitht}
  Smith, Tom, Atkinson, N.~C., native of Pender County, N.~C.

\bibitem[\protect\citeauthoryear{Shumpert, S.~B.}{}]{shumpertsb}
  Shumpert, Sarah Black, Wilmington, N.~C., native of Pender County, N.~C.

\bibitem[\protect\citeauthoryear{Stevens, J.~G.}{}]{stevensjg}
  Stevens, J.~G.

\bibitem[\protect\citeauthoryear{Sykes, D.~M.}{}]{sykesdm}
  Sykes, Dorothy Moore, Ivanhoe, N.~C., native of Sampson County, N.~C.

\bibitem[\protect\citeauthoryear{Ward, M.~B.~Jr.}{}]{wardmb}
  Ward, M.~B.~Jr., Wilmington, N.~C., native of Pender County, N.~C.

\bibitem[\protect\citeauthoryear{Weaver, E.}{}]{weavere}
  Weaver, Edward, Wilmington, N.~C., New Hanover County

\bibitem[\protect\citeauthoryear{Wright, L.}{}]{wrightl}
  Wright, Lillie, Atkinson, N.~C., native of Pender County, N.~C.

\end{thebibliography}

\printindex{}

\end{document}
